\documentclass[11pt]{article}
\usepackage[margin=1in]{geometry}
\usepackage{amsmath,amssymb}

\title{Notes on MATH 375: Probability Theory}
\author{Mariana Restrepo}
\date{Fall 2026}

\begin{document}
\maketitle

\begin{center}
\large Lecture 2
\end{center}

\tableofcontents

\section{Lecture 2: Counting and Probability}

\subsection{Permutations with repeated objects}

If $n$ objects include groups of identical objects, arranging all
$n$ positions gives $n!$ orderings, but swapping identical objects
does not create a new arrangement. Divide by the factorial of each
group size.

\paragraph{Example: SUCCESS.}
The word SUCCESS has seven letters: three S's, two C's, one U, and
one E. The number of distinct arrangements is
\[
\frac{7!}{3!\,2!\,1!\,1!}=420.
\]

\subsection{Combinations}

A \textbf{combination} selects objects without regard to their order.
The number of ways to choose $r$ objects from $n$ distinct objects is
\[
\binom{n}{r}
=\frac{n!}{r!(n-r)!}.
\]
Here $n$ is the number available and $r$ is the number chosen.

For example,
\[
\binom{5}{3}=\frac{5!}{3!\,2!}=10.
\]
Choosing the three objects to include is equivalent to choosing the
two objects to leave out, so
\[
\binom{n}{r}=\binom{n}{n-r}.
\]

\paragraph{Example: a required letter.}
How many ways can we select eight letters from the English alphabet
if A must be included? A is already selected, so choose the other
seven letters from the remaining 25:
\[
\boxed{\binom{25}{7}}.
\]

\paragraph{Example: exactly three zeros.}
A bit string of length 10 with exactly three zeros is determined by
the three positions containing zeros:
\[
\binom{10}{3}=120.
\]

\paragraph{Example: at least three zeros.}
Start with all bit strings of length 10. Exclude those with zero,
one, or two zeros:
\[
\begin{aligned}
\#(\text{at least three zeros})
&=2^{10}
 -\left[\binom{10}{0}+\binom{10}{1}+\binom{10}{2}\right]\\
&=1024-(1+10+45)\\
&=\boxed{968}.
\end{aligned}
\]
This is the \textbf{complement method}: count everything, then
subtract the cases that fail the condition.

\subsection{Sample spaces and events}

The \textbf{sample space} $S$ is the set of all possible outcomes.
An \textbf{event} $E$ is a subset of $S$: the outcomes that satisfy
a particular condition.

When all outcomes in a finite sample space are equally likely,
\[
P(E)=\frac{|E|}{|S|}
=\frac{\text{number of favorable outcomes}}
       {\text{number of possible outcomes}}.
\]
The denominator must count the outcomes in the sample space, not
just the outcomes in the event.

\paragraph{Example: rolling a die.}
For a fair six-sided die,
\[
S=\{1,2,3,4,5,6\}.
\]
Let $E$ be the event ``the result is odd.'' Then
\[
E=\{1,3,5\},
\qquad
P(E)=\frac{|E|}{|S|}=\frac{3}{6}=\frac12.
\]

\subsection{Basic probability rules}

For every event $E$,
\[
0\leq P(E)\leq 1.
\]

The probability of the entire sample space is
\[
P(S)=1.
\]

If $E$ and $F$ are \textbf{mutually exclusive}, they cannot occur
together: $E\cap F=\varnothing$. In that case,
\[
P(E\cup F)=P(E)+P(F).
\]

\subsection{Union, intersection, and complement}

$E\cup F$ means \textbf{E or F or both} occur.

$E\cap F$ means \textbf{E and F} occur together.

The complement $E^c$ means \textbf{E does not occur}. Since $E$
and $E^c$ are mutually exclusive and together make up $S$,
\[
P(E)+P(E^c)=1.
\]
Therefore, the \textbf{complement rule} is
\[
\boxed{P(E^c)=1-P(E)}.
\]

For example, if $P(E)=\frac16$, then
\[
P(E^c)=1-\frac16=\frac56.
\]

\end{document}